% 06_memory
\section{内存与分页}

分页是现代内核管理内存的基础机制。进程看到的是连续的虚拟地址空间，硬件的内存管理单元 MMU 把每个虚拟地址按页翻译到某个物理页帧，翻译关系保存在页表中。RK3588 上一页为 4 KiB，一次翻译要遍历一棵多级页表。为避免每次访存都遍历页表，MMU 内置 TLB 缓存最近用过的翻译。内核在进程创建时并不预先为全部虚拟地址建立映射，仅在某页首次被访问时经缺页异常按需建立，即首次触碰（first-touch）。首次触碰这条路径的每页固定开销，直接决定内存密集负载的表现：一段大内存被顺序写过一遍时，成千上万次缺页逐页累加，每页多花几百纳秒即放大为秒级差距。StarryOS 此前把 128 MB 私有匿名内存全部 fault-in 需 0.8 到 1.3 s，约为 Linux 同机同规模读数的 9 到 15 倍。现已在首次触碰路径上补齐三项优化，并在落地过程中修复一类潜伏于共享多架构 crate 的页表与 TLB 正确性缺陷。完成后 128 MB 私有匿名首次触碰为 \keyres{\nThpFT{} s，快于 Linux 的 \nThpFTLin{} s（\nThpRatio{} 倍）}，是在 RK3588 上领先 Linux 的少数结果之一。本节先说明分页、首次触碰、多级页表与 TLB 各自的机制，再说明既有路径的每页开销与新增的三项优化，最后说明一类共享 crate 的正确性修复。

\figphl{../figures/out/Fmem.png}{内存首次触碰建立映射与多级页表遍历路径；三项优化作用于每页固定开销，正确性修复位于共享多架构 crate 的页表遍历环节，青色为新增部分}{fig:mem-arch}

图~\ref{fig:mem-arch} 给出首次触碰这条路径及其上的多级页表遍历环节。一次缺页需分配物理帧、清零帧内容、写入页表项、并维护 TLB，其中新建映射的 TLB 失效、帧的清零方式与映射的页粒度均为每页固定成本。三项新增优化分别压缩这三处成本，共享多架构 crate 的页表遍历与 TLB 维护则是正确性修复所在。

\subsection{分页、首次触碰、多级页表与 TLB}

\paragraph{多级页表与地址翻译}
aarch64 在 4 KiB 页、48 位虚拟地址下使用四级页表。一个虚拟地址被切成四段 9 位的索引外加 12 位页内偏移，MMU 从页表基址寄存器出发，逐级用索引选出下一级表的物理地址，走满四级得到叶级页表项，其中的物理帧号与偏移拼出最终物理地址。中间各级的表项指向下一级表，叶级的表项则描述一页的物理帧与权限。某些级上的表项可直接描述一个更大的块而不再下探，例如二级上的一个块表项直接映射 2 MB，这即大页的硬件基础。这套遍历完全由硬件完成，但每翻译一个未命中的地址都要多次访存读表，代价可观。

\paragraph{TLB 与失效}
为省去重复走查，MMU 用 TLB 缓存虚拟页到物理帧的翻译。命中时翻译近乎免费，未命中时硬件走一遍页表并把结果填入 TLB。当内核改动一条已缓存的翻译时，必须失效 TLB 中的对应条目，否则硬件会继续使用陈旧翻译。aarch64 的失效指令 \texttt{tlbi} 有本地与广播两种形态：本地形态（如 \texttt{vmalle1}）只失效当前核的 TLB，广播形态（如 \texttt{vmalle1is}、\texttt{vaae1is}）经内部共享域通知所有核。SMP 下改动一条被多核共享的翻译必须用广播形态，否则兄弟核会遗留陈旧条目。失效指令前后还需 \texttt{dsb} 屏障以确保描述符写入与失效在正确次序上对其余核可见。

\paragraph{首次触碰与按需缺页}
内核为一段匿名内存建立虚拟区间时并不立即分配物理帧，页表项处于 not-present 状态。进程首次访问该区间某页时，MMU 因该页无映射触发缺页异常，缺页处理器分配一个物理帧、清零其内容以免泄漏残留数据、把帧写入叶级页表项使其转为 present，随后返回重执触发访问的指令。这条路径对每页各走一遍，因此其每页固定开销乘以页数即为触碰整段内存的总代价。StarryOS 已有 fault-around 机制，在一次缺页中顺带为邻近的 32 页预建映射，缺页次数因而受抑制，差距落在每页开销本身。

\paragraph{每页固定开销如何决定内存密集负载}
把 128 MB 私有匿名内存顺序写一遍要触碰约 32768 个 4 KiB 页。每页固定开销中的常量项在这个量级上被放大：一次多余的广播失效、一次逐字节清零、一次本可合并的缺页，单看只是几十至几百纳秒，乘以数万页即成秒级差距。sysbench 的默认 1K 块内存子测曾读到约 42 MiB/s、对 Linux 呈约 200 倍差距，此现象一度被疑为 DRAM 损坏。经核查，热态带宽为 7.3 GB/s，DRAM 本身并无问题，瓶颈落在首次触碰的缺页路径，其逐页开销即下文三项优化的对象。

\subsection{首次触碰路径的既有每页开销}

既有首次触碰路径的每页固定开销由三部分构成。其一，每建立一条全新映射均伴随一次广播 TLB 失效，128 MB 的连续触碰会产生约 32768 次 \texttt{tlbi vaae1is; dsb sy; isb} 的失效风暴；然而这些页刚从 not-present 转为 present，此前并不存在旧翻译可供陈旧，此次广播纯属多余开销。其二，新帧的清零采用普通存储逐字写零，未使用 aarch64 提供的块清零指令。其三，映射以 4K 粒度建立，每 2 MB 需经历 512 次独立缺页，缺页处理的入口开销随之乘以 512。三项均为既有路径的固定成本，下文三项优化逐一对应。

\subsection{三项首次触碰优化}

\paragraph{新建映射跳过广播 TLB 失效}
在共享的 \texttt{page\_table\_multiarch} 中为 \texttt{PagingMetaData} \ours{增加 \texttt{NEED\_FLUSH\_ON\_MAP} 标志}，默认为真，在 aarch64 与 x86\_64 上置为假。它只放行 not-present 到 present 的新建映射跳过失效，present 到 present 的 remap、protect 与 unmap 仍照常失效。跳过之所以安全，在于一条刚建立的映射此前不存在任何翻译进入 TLB，硬件不会命中一条尚未缓存的陈旧翻译；真正需要失效的情形是改写或撤销一条已可能被缓存的映射，这些情形未被放行。此项消除了整段首次触碰中约 32768 次逐页广播 \texttt{tlbi vaae1is; dsb sy; isb}。RISC-V 与 LoongArch 出于保守考虑保留默认的每次失效。此项将首次触碰由 0.8 到 1.3 s 压至约 0.37 s。

\paragraph{aarch64 dc zva 块清零}
新帧的清零\ours{改用 aarch64 的 \texttt{dc zva} 指令}，该指令按 \texttt{DCZID\_EL0} 报告的块大小一次清零一整块缓存行，取代逐字存储的清零循环。清零是首次触碰中不可省略的一步（须消除残留数据），\texttt{dc zva} 在不改变语义的前提下降低其每页成本。此项将首次触碰由约 0.37 s 压至约 0.26 s。此外配套接入 \texttt{madvise} 的 \texttt{MADV\_POPULATE\_WRITE}、\texttt{MADV\_POPULATE\_READ} 与 \texttt{MADV\_WILLNEED} 预取路径，令用户态可在访问前主动预建映射，把成批的首次触碰从关键路径上前移。

\paragraph{THP-lite 2 MB 私有匿名提升}
feature 门控的 THP-lite \ours{将 2 MB 对齐的私有匿名区提升为大页}，机制是用一个二级块表项直接映射整块 2 MB，取代 512 个 4 KiB 叶级表项。大页对首次触碰的收益有两重：其一，一次缺页即建立整块 2 MB 映射，将该区间的缺页次数由 512 降为 1，缺页处理的入口开销随之降为原来的 512 分之一；其二，一块 2 MB 仅由一条页表项描述，原本需 512 条，页表本身更小、走查更浅、TLB 每条目覆盖的地址范围扩大 512 倍，TLB 命中率随之提升。实现将区间按 “4K 头、2M 体、4K 尾” 切分，仅对 2 MB 对齐的体段提升，头尾以 4K 映射回退处理。它具备正确的 break-before-make 失效、拆分对 OOM 的原子性、无法提升时的 4K 回退与 COW 破坏路径，遵循 \texttt{MADV\_NOHUGEPAGE} 与 \texttt{PR\_SET\_THP\_DISABLE} 的逐区间/逐进程退出（后者跨 clone 继承），并兼顾 \texttt{MAP\_NORESERVE}。其中 \texttt{MADV\_NOHUGEPAGE} 供对大页敏感的负载按区间关闭提升，回落到 4K 路径。该实现在主机侧通过 54/54 测试。此项将首次触碰压至 \nThpFT{} s。

\begin{table}[H]\centering\small
\begin{tabular}{@{}p{8.2cm}r@{}}
\toprule
优化 & 128 MB 首次触碰 \\
\midrule
既有基线 & 0.8 到 1.3 s \\
新建映射跳过广播 TLB 失效 & 约 0.37 s \\
aarch64 \texttt{dc zva} 块清零新帧 & 约 0.26 s \\
THP-lite 2 MB 私有匿名提升 & \nThpFT{} s \\
\bottomrule
\end{tabular}
\caption{首次触碰缺页路径每页开销的逐项压缩（A76，128 MB 私有匿名）}
\end{table}

完成后的 128 MB 私有匿名首次触碰为 \nThpFT{} s，Linux 同机同频为 \nThpFTLin{} s，快 \nThpRatio{} 倍，板级多轮读数落在 0.024 到 0.033 s。收益来自硬件层面：更少的 TLB 广播、\texttt{dc zva} 块清零、每 2 MB 减少 512 次缺页；Linux 侧缺乏等价的用户态手段获得同等收益，因而这是一处真实的硬件领先。此处测量的是首次触碰延迟这一维度。多线程流式内存带宽属于另一维度，受板级固件缺少 DDR 变频接口所限，归入频率电压一节讨论。

\figph{../figures/out/F08.png}{THP 提升与新建映射跳过 TLB 失效把 128 MB 私有匿名首次触碰压至 \nThpFT{} s，快于 Linux 的 \nThpFTLin{} s（领先 \nThpRatio{} 倍）}

\subsection{页表与 TLB 的一类正确性修复}

提升大页、跳过失效、引入 2 MB 块这几项改动落地时，配套的多代理对抗审查（每项改动 7 到 26 名独立审查者）发现一批潜伏于共享多架构 crate 的 SMP 与 UAF 缺陷，其中若干项为既有代码中长期存在的问题。这批修复已按单一关注点拆分，准备为独立 PR 提交。逐处说明其机制如下。

\paragraph{not-present 巨页块误判为页表导致的 UAF}
最主要的一处位于 \texttt{page\_table\_multiarch} 与 \texttt{page\_table\_entry}。一个 not-present 的巨页块（例如一段被 \texttt{PROT\_NONE} 设置的 2 MB THP）在回读时 \texttt{is\_huge()} 返回假，页表遍历遂将该块对应的数据帧误当作下一级页表逐层下探。unmap 时 \texttt{dealloc\_tree} 与 \texttt{Drop} 会顺着这条错误的下探把该数据帧当作页表帧释放，若该帧正与另一进程 COW 共享，即构成一次 use-after-free。修复方式为\ours{新增 \texttt{GenericPTE::is\_table()}}（present 且非块方视为页表），所有遍历与释放判定统一经由该接口，并补齐 not-present 巨页块的 unmap 路径，仅释放块自身的帧而不再深入。同一轮审查一并修复了两处相邻缺陷：懒页（从未 fault）上的 \texttt{dealloc\_frame(0)} 回归，以 \texttt{paddr()==0} 守卫；以及 unmap 与巨页重提升时对空的中间/拆分页表的回收，其中一版基于 \texttt{is\_present} 的丢弃自身即是 UAF，已回退重写。

\paragraph{aarch64 的 break-before-make 与内存序}
另有几处 aarch64 侧的隐患，均处于内核范围，波及所有 ArceOS 家族内核而不限于 RK3588：

\begin{itemize}
\item PTE 写入与 \texttt{tlbi} 之间缺少 \texttt{dsb ishst}，失效可能在描述符写入落地之前被其余核观察到，破坏内核中每一处 break-before-make 序列；已将存储排序至 \texttt{tlbi} 之前。
\item 全量刷新回退使用了本地的 \texttt{vmalle1}，本应使用广播的 \texttt{vmalle1is}，一次触发全量刷新的大 \texttt{mprotect} 会在兄弟核上遗留陈旧 TLB；已改为广播。
\item THP 拆分对 OOM 不具原子性，先 unmap 再分配叶级页表，分配失败或 SMP 下 BBM 的 TLB 冲突窗口会遗留撕裂映射或泄漏；已改为先预留叶表的免分配拆分原语，提交阶段不可失败且可回滚。
\item 回收阶段的多次 TLB 失效批处理至单个屏障之后（\texttt{flush\_tlb\_nosync} 配合 \texttt{sync}），拆卸时每 GiB 的 \texttt{dsb} 由约 512 次降至约 16 次。
\end{itemize}

\paragraph{fork 路径上的 EFAULT}
内存工作还在板上暴露并修复了数处 fork/EFAULT 缺陷。其一，COW 引用计数使用 \texttt{u8}，当一个只读的 libc/text 帧被父进程外加约 254 个 fork 子进程共享时发生溢出，\texttt{clone\_map} 返回 \texttt{BadAddress} 转为 EFAULT；\ours{扩展为 \texttt{u32} 并采用 \texttt{checked\_add}}（溢出转 ENOMEM）后，板级 hackbench \texttt{-P g10}（400 进程）\keyres{由 240 s 超时恢复至 2.5 s}。其二，\texttt{populate\_area} 将调用方的 4 KiB 对齐范围直接传入 2 MiB 的 THP(COW) 区，\texttt{CowBackend::populate} 报 \texttt{InvalidInput} 被误映为 EFAULT，表现为板上偶发、随负载出现的 hackbench/schbench 失败（4 字节的 futex 字与 clone 的 tid 指针命中按需路径），将填充范围对齐至区的页大小即修复，QEMU smp8 复测后 \texttt{BadAddr=0}。此外在板上复现并修复了两例与 \texttt{PROT\_NONE} 巨页相关的路径：fork 一个曾对 THP 区做过 \texttt{PROT\_NONE} 的进程触发的 EFAULT，以及穿过一个 not-present 巨页块的部分操作触发的 EFAULT，二者均随上文 \texttt{is\_table()} 与 not-present 巨页帧回收一并收敛。整区 \texttt{PROT\_NONE} 再放开后 \texttt{mprotect} 返回 0、随后访问触发 SIGSEGV 的一例，此前一度挂起，现\ours{以整区放开时对 not-present 巨页块的重新提升修复}，thp\_reenable 的 A 至 F 六个用例在板上与 Linux 一致通过。另修复两处 SMP 下的数据完整性缺陷：匿名区在 \texttt{mprotect(+W)} 放宽写权限时未破坏 COW，可致跨进程写穿，改为在放宽写权限时对共享帧执行 COW 破坏；futex 字此前为非原子读取，SMP 下可能读到撕裂值而丢失唤醒，改以 \texttt{atomic\_read\_user\_u32} 原子读取。

这批修复均位于 ArceOS、StarryOS、Axvisor 三套系统共享的多架构 crate 中，波及四种架构，一处改动即令三套系统受益。此类页表与 TLB 的 UAF 及乱序能够通过常规浅层测试，却在 SMP 叠加 COW 与 THP 时悄然破坏内存，能够将其定位，正是原子 PR 与多代理对抗审查这一纪律的直接产物；一轮 8 代理专门排查投机取巧实现的对抗审查在最终交付集上未发现投机式实现。

\figph{../figures/out/F09.png}{页表与 TLB 修复位于共享多架构 crate，一处改动波及 ArceOS/StarryOS/Axvisor 与四种架构}
